        	\begin{tikzpicture}[
			font=\sffamily,
			line width=1pt,
			label_node/.style={fill=white, draw=black!50, rounded corners=2pt, inner sep=4pt, text=black, font=\footnotesize\bfseries},
			arrow/.style={thick, -Latex, black!70}
			]
			
			% ==========================================
			% ANSICHT: VORNE
			% ==========================================
			\begin{scope}[shift={(0,0)}]
				% Titel
				\node at (2, -1.5) {\textbf{Ansicht: Vorne}};
				
				% Gehäuse (Hauptkörper)
				\draw[fill=black!5, rounded corners=3pt] (0,0) rectangle (4,4);
				
				% Dreheinheit (auf dem Gehäuse)
				\draw[fill=black!15, thick] (1.2, 4) rectangle (2.8, 4.6);
				
				% Ultraschallsensor (auf der Dreheinheit)
				\draw[fill=black!10, thick] (0.1, 4.4) rectangle (3.9, 5.8);
				% Sensor-Augen (Sender/Empfänger)
				\draw[fill=white, thick] (1.4, 5.1) circle (0.35);
				\draw[fill=white, thick] (2.6, 5.1) circle (0.35);
				
				% --- Beschriftungen Vorne ---
				\node[label_node] (lblSensor) at (-1.5, 5.1) {Ultraschallsensor};
				\draw[arrow] (lblSensor) -- (0.8, 5.1);
				
				\node[label_node] (lblDreh) at (-1.6, 4.3) {Dreheinheit};
				\draw[arrow] (lblDreh) -- (1.1, 4.3);
				
				\node[label_node] (lblGeh) at (-1.5, 2) {Geh\"ause};
				\draw[arrow] (lblGeh) -- (0, 2);
			\end{scope}
			
			% ==========================================
			% ANSICHT: HINTEN
			% ==========================================
			\begin{scope}[shift={(5,0)}]
				% Titel
				\node at (2, -1.5) {\textbf{Ansicht: Hinten}};
				
				% Gehäuse (Hauptkörper)
				\draw[fill=black!5, rounded corners=3pt] (0,0) rectangle (4,4);
				
				% Öffnungen Hinten (LEDs und USB-C)
				% USB-C Port (unten mittig)
				\draw[fill=white, thick, rounded corners=1pt] (1.6, 0.6) rectangle (2.4, 0.9);
				\draw[fill=black!70] (1.7, 0.7) rectangle (2.3, 0.8); % Kontaktzunge
				
				% LEDs (oben)
				\draw[fill=white, thick] (2.5, 3) circle (0.2); % LED rot Position
				\draw[fill=white, thick] (3, 3) circle (0.2); % LED grün Position
				
				% Dreheinheit (Rückseite)
				\draw[fill=black!15, thick] (1.2, 4) rectangle (2.8, 4.6);
				
				% Ultraschallsensor (Rückseite - keine Augen sichtbar)
				\draw[fill=black!10, thick] (0.1, 4.4) rectangle (3.9, 5.8);
				% Rückseite des PCBs (kleine Details)
				%\draw[fill=black!20, draw=none] (1.2, 4.8) rectangle (1.6, 5.4);
				%\draw[fill=black!20, draw=none] (2.4, 4.8) rectangle (2.8, 5.4);
				
				% --- Beschriftungen Hinten ---
				\node[label_node] (lblRed) at (0.5, 3) {LED rot};
				\draw[arrow] (lblRed) -- (2.3, 3);
				
				\node[label_node] (lblGreen) at (5, 3) {LED gr\"un};
				\draw[arrow] (lblGreen) -- (3.2, 3);
				
				\node[label_node] (lblUSB) at (5, 0.75) {USB-C Port};
				\draw[arrow] (lblUSB) -- (2.6, 0.75);
			\end{scope}
			
		\end{tikzpicture}%